\documentclass[10pt]{amsart}

\usepackage[margin=1in]{geometry}
\usepackage[T1]{fontenc}
\usepackage{lmodern}
\usepackage{amsmath,amssymb,amsthm,mathtools}
\usepackage{microtype}
\usepackage{url}
\usepackage{enumitem}
\usepackage{booktabs}
\usepackage{array}
\usepackage[dvipsnames]{xcolor}
\usepackage[colorlinks=true,linkcolor=MidnightBlue,citecolor=OliveGreen,urlcolor=MidnightBlue]{hyperref}
\usepackage[capitalize,nameinlink]{cleveref}

\setlist{leftmargin=*,itemsep=2pt,topsep=4pt}
\emergencystretch=3em
\raggedbottom
\Urlmuskip=0mu plus 1mu

% Theorem environments; the optional note (status) is set at the start of the body.
\theoremstyle{plain}
\newtheorem{thmbase}{Theorem}[section]
\newtheorem{theoremx}[thmbase]{Theorem}
\newtheorem{propositionx}[thmbase]{Proposition}
\newtheorem{lemmax}[thmbase]{Lemma}
\newtheorem{corollaryx}[thmbase]{Corollary}
\newtheorem{factx}[thmbase]{Fact}
\newtheorem{importedx}[thmbase]{Imported statement}
\theoremstyle{definition}
\newtheorem{definitionx}[thmbase]{Definition}
\theoremstyle{remark}
\newtheorem{remarkx}[thmbase]{Remark}
\newcommand{\thmnoteinbody}[1]{\if\relax\detokenize{#1}\relax\else{\upshape[#1]}\enspace\fi}
\newcommand{\defthm}[1]{%
  \newenvironment{#1}[1][]{\begin{#1x}\thmnoteinbody{##1}\ignorespaces}{\end{#1x}}}
\defthm{theorem}\defthm{proposition}\defthm{lemma}\defthm{corollary}\defthm{fact}
\defthm{imported}\defthm{definition}\defthm{remark}
\crefname{theoremx}{Theorem}{Theorems}
\crefname{propositionx}{Proposition}{Propositions}
\crefname{lemmax}{Lemma}{Lemmas}
\crefname{corollaryx}{Corollary}{Corollaries}
\crefname{factx}{Fact}{Facts}
\crefname{importedx}{Imported statement}{Imported statements}
\crefname{definitionx}{Definition}{Definitions}
\crefname{remarkx}{Remark}{Remarks}
\numberwithin{equation}{section}

\newcommand{\F}{\mathbb{F}}
\newcommand{\R}{\mathbb{R}}
\newcommand{\Z}{\mathbb{Z}}
\newcommand{\Q}{\mathbb{Q}}
\newcommand{\cC}{\mathcal{C}}
\newcommand{\cN}{\mathcal{N}}
\newcommand{\cF}{\mathfrak{F}}
\newcommand{\rk}{\operatorname{rk}}
\newcommand{\Mat}{\operatorname{Mat}}
\newcommand{\GL}{\operatorname{GL}}
\newcommand{\cost}{\operatorname{cost}}
\newcommand{\Fcat}{F_{\mathrm{CAT}}}
\newcommand{\eps}{\varepsilon}
\newcommand{\dd}[2]{|#1\,#2|}
% status labels
\newcommand{\stP}{proved here}
\newcommand{\stPL}{proved here and in Lean}
\newcommand{\stR}{refereed by an AI referee only}
\newcommand{\stN}{numerical}
\newcommand{\stU}{not refereed}
\newcommand{\cond}[1]{conditional on #1}
% sources in the repository
\DeclareUrlCommand\src{\def\UrlFont{\small\ttfamily}\def\UrlBreaks{\do\/\do\-\do\.}\def\UrlBigBreaks{}}
\DeclareUrlCommand\lean{\urlstyle{tt}\def\UrlBreaks{\do\.\do\_}\def\UrlBigBreaks{}}
% repository commit
\newcommand{\commit}{e97d95610af61382c28171bda4d9ebeb18345051}
\edef\repourl{https://github.com/xennygrimmato/math/tree/\commit/research/sat3}

\begin{document}

\title[Three barriers at exponent one half]{Three barriers at exponent one half\\ for the counting engines}

\author{Vaibhav Tulsyan}

\date{27 November 2026}

\subjclass[2020]{68Q17, 68Q25, 15A69, 51E24, 53C23}
\keywords{Exponential Time Hypothesis, \#3-SAT, tensor networks, asymptotic spectrum, CAT(0) spaces, spherical buildings, nets, Lean~4}

\begin{abstract}
In October 2026 Anthropic and OpenAI each claimed a deterministic algorithm that counts the
satisfying assignments of a 3-CNF on $n$ variables in time $2^{O(n/\sqrt{\log n})}$. Anthropic's
algorithm is a recursion on clique sums whose levels merge the nodes in one of two ways, and the
exponent $1/2$ comes from a level factor $1/\sqrt2$. Levels built on $p$-nets of $\F_q^m$ have the
same node growth and $p$ merges, so they could give a smaller factor. We state four results that
explain why this does not happen with the present methods, each with its hypotheses and its status.
(1) The space of real filtrations of $\F_q^N$, a CAT(0) space, carries a placement of the $p$-net
of $\F_q^m$ ($m\ge2$, $g=q^m$) of value at least $1/\sqrt2-O(1/g)$, and a two-bimedian inequality
shows that no placement in any CAT(0) space has value above $\sqrt{g/(2(g+1))}$. (2) If Anthropic's
contracting hypothesis fails over a field $k$, every finite subset of this space over $\F_{2^s}$ is
a price over $k$; hence every identity over $k$ for a net level with $2^{o(\lambda)}$ terms has a
piece whose true cost is at least $1/\sqrt2-c/g$ per bit under every merge at once for large $\lambda$, and an identity
for one fixed eight-wire configuration that stays below $1/\sqrt2$ by a constant would prove the
contracting hypothesis. (3) For covers generated from crossing identities by nesting, at every
depth, a Steiner crossing inequality valid in all CAT(0) spaces shows that the certified cost of
some term is at least the value of every CAT(0) placement; this is a barrier for a proof method,
not for the tensors. (4) Conditionally on two unrefereed theorems of OpenAI, a straight-line program
for the clique polynomial whose exponent beats $1/2$ in the long-wire regime gives ports in
characteristic $0$ and a power-saving algorithm for 3-SAT. Cyclic nets, the difference between true
and certified costs, and non-algebraic algorithms remain open; we report the numerical evidence.
The mathematics, proofs, code and text were produced by AI systems, the referees were AI agents,
and no person has checked the proofs.
\end{abstract}

\maketitle

% =====================================================================
\section{Introduction}\label{sec:intro}
% =====================================================================

\subsection{The exponent one half}
Anthropic~\cite{AnthropicETH} and OpenAI~\cite{OpenAIETH} each claimed in October 2026 a
deterministic algorithm for \#3-SAT on $n$ variables in time $2^{O(n/\sqrt{\log n})}$. Both claims
are unrefereed. Anthropic's algorithm reduces the count to clique sums: $K$ nodes, $\lambda$ bits
on every edge, an arbitrary table at every node. One level of its recursion partitions the nodes of
a clique on $K'$ nodes into $K$ groups in one of two ways (rows or columns of a $g\times g$ grid), and
writes the product of the edge wires as a sum of $2^{o(\lambda)}$ terms, each cheap under one of the
two merges. The cost of a level, per edge bit, is a factor $\rho$, and the node count grows as
$K'=K^{\beta}$ with $\beta=2$. A tower of such levels gives an exponent $e_K=(\log K)^{-\alpha'}$
with
\begin{equation}\label{eq:alpha}
\alpha'=\log_\beta(1/\rho^*),\qquad \rho^*=\liminf\text{ of the level factors}.
\end{equation}
Anthropic's crossing identity (its Proposition~B.2 in~\cite{AnthropicVP}, Appendix~B
in~\cite{AnthropicETH}) gives $\rho=1/\sqrt2$, hence $\alpha'=1/2$ and the running time
$2^{n/(\log n)^{1/2+o(1)}}$.

A design with $p$ merges and node growth $\beta$ has, under Euclidean prices, the flat floor
$\rho\ge1/\sqrt p$ on balanced configurations, which would allow $\alpha'$ up to
$\tfrac12\log_\beta p$. For $p\le\beta$ this is at most $1/2$. The $p$-nets break this count: the
$p$ parallel classes of an affine plane of order $g$ give $p$ merges of $g^2$ nodes into $g$ groups,
so $\beta=2$ for every $p$, and the flat floor alone would allow $\alpha'=\tfrac12\log_2 3\approx0.79$
for three merges. The question is the true floor $\rho^*$ of net levels.

\subsection{Results}
We consolidate one day's results on this question, refereed by AI agents, into four statements,
ordered by strength of hypothesis. Here $g=q^m$ is the order of the net and the value of a placement
is the lower bound per edge bit that the price method extracts from it (\cref{sec:setting}).

\begin{enumerate}[label=(\arabic*)]
\item \emph{Theorem A and the cap (\cref{sec:thmA}).} For every prime power $q$, every $m\ge2$ and
every set of $p\ge2$ slopes, the $p$-net of $\F_q^m$ has a placement in the space of real
filtrations of $\F_q^{m+2}$, a CAT(0) space, of value at least
$1/\sqrt2-((q+1)/\sqrt2-1)/(g+1)$. No placement in any CAT(0) space has value above
$\sqrt{g/(2(g+1))}$. On every balanced set of generic edges both bounds equal $1/\sqrt2$. No
hypothesis is used.
\item \emph{The price barrier (\cref{sec:price}).} If the contracting hypothesis of~\cite{AnthropicVP}
fails over $k$, every finite subset of the filtration space over $\F_{2^s}$ is a price over $k$. On
nets of $\F_2^m$ every identity over $k$ for the level tensor with $2^{o(\lambda)}$ terms then has a
piece of true cost at least $(1/\sqrt2-1.13/(g+1)-\eps)\lambda$ per edge, under all three merges at
once, for $\lambda\ge\lambda_0(g,\eps)$. Contrapositively, identities for one fixed configuration of
eight wires ($E_D$) whose pieces all stay below $(1/\sqrt2-\eta)8\lambda$ under some merge would
prove the contracting hypothesis, hence a power-saving algorithm for 3-SAT.
\item \emph{Theorem W (\cref{sec:thmW}).} For every placement in a CAT(0) space and every cover
generated from circle-covering crossing identities by routing and nesting, some term has certified
cost at least the placement's bound under every merge at once. With the zero-loss lemma
(\cref{lem:zeroloss}, refereed after the other results) this holds exactly at all depths. Certified costs are upper
bounds on true costs, so Theorem~W is a barrier for a proof method.
\item \emph{Theorem T (\cref{sec:thmT}).} If OpenAI's Theorem~3.1 of~\cite{OpenAIHC} holds as
printed, every straight-line program computing the clique polynomial $f_{K,\lambda}$ as a polynomial
identity, over a field whose characteristic has no ports, needs $2^{cK^2\lambda/\sqrt{\log K}}$ gates
for large even $K$ and $\lambda\ge CK(\log K)^{3/2}$. With OpenAI's Theorem~7.1 of~\cite{OpenAIIM}, such
an algorithm over $\Q$ with $e(K)\sqrt{\log K}\to0$ would give 3-SAT in time $2^{O(v^{1-\delta})}$.
\end{enumerate}

Statements (1) and (2) concern the price method, which bounds true costs; (1) bounds what any price
can prove and supplies the extremal geometry, and (2) turns that geometry into a lower bound when
contracting fails. Statement (3) needs no hypothesis but bounds only certified costs. Statement (4)
covers all algebraic algorithms in its regime, but rests on unrefereed work of another group.
\Cref{sec:open} lists what is open, and \cref{sec:limits} the limitations.

\subsection{Status of each claim}\label{sec:status}
Every statement carries one of the following labels.
\emph{Proved here}: a complete proof is in this paper; unless stated otherwise it was also checked by
an AI referee. \emph{Proved here and in Lean}: in addition, the part named is formalised in our
Lean~4 library~\cite{Lean4,Mathlib20}, with the scope stated in \cref{sec:lean}. \emph{Refereed by an
AI referee only}: the full proof is in the repository notes and was checked by an AI referee; here we
give a sketch with pointers. \emph{Conditional on $X$}: the statement assumes $X$. \emph{Numerical}:
a computation, with the script named in the repository. \emph{Not refereed}: written up but not yet
checked by a referee. No statement of this paper has been checked by a person. The repository files
are cited as \src{q1-notes.md}~\S2 and so on, relative to \src{research/sat3/} unless a path from the repository root
is given. The file \src{paper-barriers/CHECKS.md} lists every statement with its source and status.

Several inputs are unrefereed papers of the two labs. The statements we use from
them are quoted and marked as imported. Anthropic's paper~\cite{AnthropicVP} states that its proofs
were found and written by AI and that no human mathematician has checked them; OpenAI's release
describes its manuscripts as raw research data.

% =====================================================================
\section{Setting}\label{sec:setting}
% =====================================================================

\subsection{Configurations, merges and costs}
A \emph{configuration} is a finite set of nodes with a set $E$ of edges (unordered pairs of nodes).
Each edge $e=\{n,n'\}$ carries $\lambda$ bits and has two \emph{legs}, $e^+$ at $n$ and $e^-$ at
$n'$. Its \emph{level tensor} over a field $k$ is $f=\prod_{e}\langle2^\lambda\rangle_{e^+e^-}$, the
product of identity matrices (wires) on the legs. A \emph{merge} is a partition of the nodes into
groups. For a tensor $A$ on the legs and a merge $k$, the \emph{D1 cost} $c_k(A)$ is the least
$\sum_{G<G'}\log_2 d_{GG'}$ such that the push-forward of $A$ to the groups (a group holds the
indices of all legs of its nodes) is a restriction of the graph tensor
$\prod_{G<G'}\langle d_{GG'}\rangle_{GG'}$; it is the cost of the merged clique instance. A
\emph{piece family} is an identity $\sum_{i\in I}A_i=\zeta f$ over $k$ with $\zeta\ne0$. Its
\emph{level value} is $\max_i\min_k c_k(A_i)$ and $\rho$ is the level value divided by $\lambda|E|$.

An \emph{engine} in the sense of \src{research/ports/c1d1/notes.md}~D1.1 is, for each $n$, a tower of levels
with node counts $K_{\ell+1}\ge K_\ell^\beta$ and level factors $\rho_\ell$. It runs the recursion
$e_{\ell+1}=\rho_\ell e_\ell+(\text{non-negative terms})$, and its \emph{log-exponent} is
$\alpha'=\liminf_n(-\log e_L)/\log\log K_L$. We call $e_L$ the \emph{recursion exponent}; it is
computed from the true costs $c_k$. In a level of class $(p,\beta,\mathcal C)$ the edges are partitioned
into copies of members of a finite list $\mathcal C$ of configurations, each copy with every edge
crossing every merge, plus a remainder carrying a fraction $\xi$ of the bits; for each $C\in\mathcal C$
there is one identity with at most $2^{\tau(\lambda)}$ pieces, $\tau(\lambda)=o(\lambda)$, and the terms
of the level are the products of one piece per copy.

\subsection{Nets}\label{sec:nets}
Let $q$ be a prime power, $m\ge2$, $g=q^m$, and $P\subseteq\mathbb P^1(\F_q)$ a set of $p\ge2$ slopes,
with a representative $x\in\F_q^2$ for each. The \emph{$p$-net} $\cN=\cN(q,m,P)$ has the $g^2$ nodes
$a\in\Mat_{m\times2}(\F_q)$; its merge of class $x$ is the partition into the $g$ lines
$\{a:ax=y\}$, $y\in\F_q^m$. Its level is the complete graph on the nodes,
$E=g^2(g^2-1)/2$. Two nodes $a,a'$ form a \emph{generic} edge if $\rk(a-a')=2$; then they lie on no
common line, and the edge crosses all $q+1$ classes. For $q=2$ and $P=\{0,\infty,1\}$ this is the
3-net of the group $\F_2^m$, with classes rows, columns and their sum. Counting matrices of rank
$2$ and $1$,
\begin{equation}\label{eq:counts}
E_{\mathrm{gen}}=E\,\frac{g-q}{g+1},\qquad E_{\mathrm{coll}}=E\,\frac{p}{g+1},\qquad
E_{\mathrm{other}}=E\,\frac{q+1-p}{g+1},
\end{equation}
for generic edges, edges on a common line, and the remaining rank-one edges. The configuration
$X_q(P)$ is the set of generic edges of $\cN(q,2,P)$. For $q=2$ it is $X_4$: $16$ nodes $\Mat_2(\F_2)$,
$48$ edges $\{a,a'\}$ with $a-a'\in\GL_2(\F_2)$, three merges into four groups of four. For
$D\in\GL_2(\F_2)$ the configuration $E_D=\{\{a,a+D\}\}\subset X_4$ is a perfect matching: eight
disjoint wires, each crossing all three merges.

\subsection{Placements, savings and balance}
A \emph{placement} of a configuration in a metric space $Y$ assigns a point to every leg and a hub
$o^k_G$ to every group $G$ of every merge $k$; write $o^k(n)$ for the hub of the group of $n$. Put
\[
L=\sum_e\dd{e^+}{e^-},\quad H_k=\sum_{\text{legs }\ell}\dd{\ell}{o^k(\ell)},\quad
\delta_k=\max_{G\ne G'}\dd{o^k_G}{o^k_{G'}},\quad b_k=\frac{L-H_k}{\delta_k},
\]
and $s^k_e=\dd{e^+}{e^-}-\dd{e^+}{o^k(n)}-\dd{e^-}{o^k(n')}$, the saving of $e$ under $k$, so that
$b_k\delta_k=\sum_es^k_e$. The \emph{value} of the placement is $\min_kb_k/|E|$, and $\Fcat$ of a
configuration is the supremum of the value over placements in CAT(0) spaces.

\begin{definition}\label{def:balance}
A configuration is \emph{balanced} if for all groups $G,G'$ of two different merges the number of
ordered pairs $(y,z)$ with $\{y,z\}\in E$, $y\in G$, $z\in G'$ equals $\sum_{n\in G\cap G'}\deg n$.
\end{definition}

The full level of a $p$-net is balanced for every pair of merges: two lines of different classes meet
in one node $n_0$, and both sides equal $g^2-1$. A single wire is not balanced.

\subsection{Spectral points, prices and the two hypotheses}\label{sec:prices}
We use the definitions of~\cite{AnthropicVP}, Appendix~A. A \emph{spectral point} $\varphi$ on a finite
set $V$ of parties is a function on $V$-tensors that is multiplicative, takes the value $1$ on the
scalar $1$, and is subadditive on covers. On wires $\varphi(\langle q\rangle_{uv})=q^{d_\varphi(u,v)}$
for a pseudometric $d_\varphi$ with values in $[0,1]$ (its Lemma~A.2). A \emph{price} on $V$ is a
pointwise limit of functions $c\,d_\varphi$ with $c\ge0$ (its Definition~A.4); a function on the pairs of
an arbitrary set is a price if all its pull-backs to finite sets are. Pull-backs, non-negative
multiples and limits of prices are prices. A CAT(0) space is a complete metric space in which any two
points $x,v$ have a midpoint satisfying the CN inequality; then (its Lemma~C.4(a), and~\cite{BH99})
for every $t\in[0,1]$ there is a unique point $\gamma_t(x,v)$ at distance $t\dd xv$ from $x$ and
$(1-t)\dd xv$ from $v$, and for every $z$
\begin{equation}\label{eq:CN}
\dd z{\gamma_t(x,v)}^2\le(1-t)\dd zx^2+t\dd zv^2-t(1-t)\dd xv^2 .
\end{equation}
A \emph{price space} is a CAT(0) space whose metric is a price.

\begin{imported}[Anthropic~\cite{AnthropicVP}; unrefereed source]\label{imp:AN}
(a) (Prop.~C.3) Every price on a set is the pull-back of the metric of a price space.
(b) (Def.~D.1) The \emph{flat hypothesis} over $k$: all finite subsets of Euclidean spaces are prices.
The \emph{contracting hypothesis} over $k$: there are $\mu\ge2$ and $1\le s<\sqrt\mu$ with
$D(\mu)\le sD(1)$ for every price on $\Mat_\mu(\F_2)$, where $D(r)^2$ is the mean of $\dd a{a+W}^2$ over
all $a$ and all $W$ of rank $r$.
(c) (Thm.~D.2, Lean \lean{thm_alt}) Over every field one of the two hypotheses holds.
(d) (Proof of Thm.~D.2, Lean \lean{thm_alt_rank_price}) If the contracting hypothesis fails over $k$,
then $\sqrt{\rk(a-b)}$ is a price over $k$ on $\Mat_\mu(\F_2)$ for every $\mu\ge3$.
\end{imported}

The contracting hypothesis has an algorithmic consequence. By Theorem~R of
\src{research/ports/note-contracting-eth.md}~\S3 (refereed by an AI referee in its form with advice; the
uniform form rests on an audit of Anthropic's Lean proof terms by one reader), contracting over a
field of characteristic $0$ gives a deterministic algorithm for 3-SAT in time
$2^{O(n^{1-\delta})}\mathrm{poly}$; in characteristic $p$ the same note gives a randomised one
(Theorem~C1-AN there). We call such algorithms \emph{power-saving}.

\begin{proposition}[price method; \stR, \src{research/vpvnp/b2/notebook.md}~\S1, \src{notes.md}~\S2]\label{prop:price}
Let $d$ be a price over $k$ on the legs and hubs of a placement of a configuration. For every
identity $\sum_{i\in I}A_i=\zeta f$ over $k$ with $|I|=2^{o(\lambda)}$ some $i$ has
$c_k(A_i)\ge\lambda b_k-o(\lambda)$ for every merge $k$ at once; the $o(\lambda)$ depends on the
configuration and the price.
\end{proposition}

\begin{proof}[Sketch]
Approximate $d$ by $c\,d_\varphi$. Subadditivity and $\varphi(\langle2^\lambda\rangle_{uv})=2^{\lambda
d_\varphi(u,v)}$ give some $i$ with $\log\varphi(A_i)\ge\lambda\sum_ed_\varphi(e)-\log|I|$. A network of
cost $c_k(A_i)$ on the groups of $k$ routes $A_i$ through one hub per group (Lemma~A.3
of~\cite{AnthropicVP}), so $\log\varphi(A_i)\le\lambda\sum_\ell d_\varphi(\ell,o^k(\ell))+\delta_k(\varphi)
c_k(A_i)$. Multiply by $c$ and compare. The scale $c$ depends only on the finite point set.
\end{proof}

% =====================================================================
\section{Two CAT(0) inequalities and the cap}\label{sec:cap}
% =====================================================================

All prices sit in price spaces (\cref{imp:AN}(a)), so an inequality valid in every CAT(0) space bounds
what the price method can ever prove. Besides~\eqref{eq:CN} we use that $t\mapsto\dd{\gamma_t(x,x')}
{\gamma_t(y,y')}$ is convex for all $x,x',y,y'$ (\cite{BH99}, Prop.~II.2.2); in particular
$\dd z{\gamma_t(x,v)}\le(1-t)\dd zx+t\dd zv$.

\begin{proposition}[two bimedians; \stPL]\label{prop:bimedian}
In a CAT(0) space, for points $o,w,o',w'$ and $t\in[0,1]$,
\[
\dd{\gamma_t(o,w)}{\gamma_t(o',w')}^2+\dd{\gamma_t(o,w')}{\gamma_t(o',w)}^2\le2(1-t)^2\dd o{o'}^2+2t^2\dd
w{w'}^2 .
\]
\end{proposition}

\begin{proof}
Put $m=\gamma_t(o,w)$, $m'=\gamma_t(o',w')$. By~\eqref{eq:CN} for $m'$ against $m$, and then for $m$
against $o'$ and against $w'$,
\begin{align*}
\dd m{m'}^2&\le(1-t)\dd m{o'}^2+t\dd m{w'}^2-t(1-t)\dd{o'}{w'}^2\\
&\le(1-t)^2\dd o{o'}^2+t^2\dd w{w'}^2+t(1-t)\bigl[\dd w{o'}^2+\dd o{w'}^2-\dd ow^2-\dd{o'}{w'}^2\bigr].
\end{align*}
The same computation with $w$ and $w'$ exchanged bounds the second squared distance by
$(1-t)^2\dd o{o'}^2+t^2\dd w{w'}^2+t(1-t)[\dd{w'}{o'}^2+\dd ow^2-\dd o{w'}^2-\dd{o'}w^2]$. The
brackets cancel in the sum.
\end{proof}

In a Hilbert space the inequality is an equality (the parallelogram law); \cref{sec:lean} lists the
Lean statements, \lean{bimedian_weighted} and \lean{bimedian_inner_eq}.

\begin{theorem}[two-merge cap; \stP]\label{thm:cap2}
Let a configuration with two merges be balanced, and let $E_k$ be the number of edges crossing merge
$k$. For every placement in a CAT(0) space, $b_1+b_2\le\sqrt{E(E_1+E_2)}$. In particular a balanced,
fully crossing configuration has $\min(b_1,b_2)\le E/\sqrt2$, and no price over any field proves more.
\end{theorem}

\begin{proof}
Put $\lambda_k=1/\delta_k$, $\Lambda=\lambda_1+\lambda_2$, $t=\lambda_2/\Lambda$, $\mu=(1-t,t)$ and
$\beta(n)=\gamma_t(o^1(n),o^2(n))$. For an edge with legs $x$ at $n$ and $y$ at $n'$, convexity gives
$\lambda_1\dd x{o^1(n)}+\lambda_2\dd x{o^2(n)}\ge\Lambda\dd x{\beta(n)}$, so
\[
\textstyle\sum_k\lambda_ks^k_e\le\Lambda\bigl(\dd xy-\dd x{\beta(n)}-\dd y{\beta(n')}\bigr)\le\Lambda\dd
{\beta(n)}{\beta(n')}.
\]
Two applications of~\eqref{eq:CN} give
$\dd{\beta(n)}{\beta(n')}^2\le\sum_{k,l}\mu_k\mu_l\dd{o^k(n)}{o^l(n')}^2-t(1-t)\bigl(\dd{o^1(n)}{o^2(n)}^2+
\dd{o^1(n')}{o^2(n')}^2\bigr)$. Sum over the edges. The terms with $k\ne l$ add up to
$\mu_1\mu_2\sum_{G,G'}\dd{o_G}{o_{G'}}^2\cdot\#\{(y,z)\}$ and the subtracted terms to
$\mu_1\mu_2\sum_{G,G'}\dd{o_G}{o_{G'}}^2\sum_{n\in G\cap G'}\deg n$; they cancel by balance. Hence
$\sum_e\dd{\beta(n)}{\beta(n')}^2\le\sum_k\mu_k^2\sum_e\dd{o^k(n)}{o^k(n')}^2\le\sum_k\mu_k^2E_k\delta_k^2
=(E_1+E_2)/\Lambda^2$, and by Cauchy--Schwarz
$b_1+b_2=\sum_e\sum_k\lambda_ks^k_e\le\Lambda\sum_e\dd{\beta(n)}{\beta(n')}\le\sqrt{E(E_1+E_2)}$.
\end{proof}

\begin{corollary}[the cap on $p$-nets; \stP]\label{cor:cap}
For every placement of the full level of a $p$-net of order $g$ in a CAT(0) space,
$\tfrac1p\sum_kb_k\le\sqrt{E\,E_\times/2}$ with $E_\times=Eg/(g+1)$ the number of edges crossing one
merge. So every price, under any hypothesis, gives value at most $\sqrt{g/(2(g+1))}<1/\sqrt2$.
\end{corollary}

\begin{proof}
Every pair of merges is balanced (\cref{sec:nets}), so \cref{thm:cap2} gives
$b_k+b_l\le\sqrt{2EE_\times}$. Sum over the $\binom p2$ pairs. The count $E_\times$ follows
from~\eqref{eq:counts}: the edges internal to one merge are its $E/(g+1)$ collinear edges.
\end{proof}

For $K_4$ with its three perfect matchings the same argument, applied to each of the three crossings
inside it, gives $\min_kb_k\le2\sqrt2$, which the flat placement attains (\src{notes.md}
Cor.~3.3, \stR). For the configuration $\Delta_3$ (the nine edges of the fourth parallel class of $\mathrm{AG}(2,3)$
with rows, columns and diagonals as merges) it gives $\min_kb_k\le9/\sqrt2$.

% =====================================================================
\section{Theorem A: the filtration space reaches the cap}\label{sec:thmA}
% =====================================================================

\subsection{The space of filtrations}
Let $V=\F_q^N$. A \emph{filtration} is a decreasing, left-continuous family $(F_s)_{s\in\R}$ of
subspaces with $F_s=V$ for $s\ll0$ and $F_s=0$ for $s\gg0$. For a basis $f=(f_1,\dots,f_N)$ and
$x\in\R^N$ put $F^{f,x}_s=\operatorname{span}\{f_i:x_i\ge s\}$; $f$ is \emph{adapted} to $F$ if
$F=F^{f,x}$ for some $x$, the level vector.

\begin{lemma}[\stR; standard, \src{q1-notes.md}~\S1]\label{lem:filt}
(a) Any two filtrations have a common adapted basis. (b) If $f$ is adapted to $F$ and $F'$ with level
vectors $x,x'$, the number of $i$ with $x_i\ge s$ and $x'_i\ge s'$ is $\dim(F_s\cap F'_{s'})$; hence
$d(F,F')=|x-x'|$ does not depend on $f$.
\end{lemma}

Part~(a) is the existence of a common adapted basis for two flags (Bruhat decomposition); (b) holds
because $F_s\cap F'_{s'}$ is a coordinate subspace of $f$. Let $\cC_N(q)$ be the set of filtrations
with the metric $d$. A subspace $U$ is the filtration ``$U$ at level $1$, $V$ at level $0$'', and for
subspaces $d(U,U')^2=\dim(U+U')-\dim(U\cap U')$.

\begin{fact}[classical; identification checked by an AI referee, \src{review-q1.md}~\S1]\label{fact:cat}
$(\cC_N(q),d)$ is isometric to $\mathrm{Cone}(\Delta)\times\R$, where $\Delta$ is the spherical building
of $\GL_N(\F_q)$ with its round metric; the apartments are the sets $\{F^{f,x}:x\in\R^N\}$ with the
Euclidean metric, and the $\R$-factor is the mean level. Hence $\cC_N(q)$ is CAT(0): $\Delta$ is
CAT(1)~\cite[Thm.~II.10A.4]{BH99}, Euclidean cones over CAT(1) spaces are CAT(0)~\cite[Thm.~II.3.14]{BH99},
and products of CAT(0) spaces are CAT(0).
\end{fact}

\begin{lemma}[far legs at point hubs; \stP]\label{lem:farleg}
Let $f$ be a basis, $u\in\R^N$ a unit vector, $c\in\R^N$, $x_T=F^{f,c+Tu}$ and $y_T=F^{f,c-Tu}$. Let
$h=\langle\sum_{i\in A}\alpha_if_i\rangle$ with all $\alpha_i\ne0$ be a point, taken as the filtration
``$h$ at level $1$, $V$ at level $0$''. Then $d(x_T,h)=T+\langle c,u\rangle-\min_{i\in A}u_i+O(1/T)$ and
$d(y_T,h)=T-\langle c,u\rangle+\max_{i\in A}u_i+O(1/T)$.
\end{lemma}

\begin{proof}
Let $i^*\in A$ minimise $c_i+Tu_i$ over $A$; for large $T$ it minimises $u_i$. Replace $f_{i^*}$ by
$g=\sum_{A}\alpha_if_i$. For $s\le c_{i^*}+Tu_{i^*}$ all $f_i$, $i\in A$, lie in $F_s$, so $g$ and the
old $f_{i^*}$ lie in the span of the new basis vectors of level $\ge s$; for larger $s$ neither does.
So the new basis is adapted to $x_T$ with the same level vector, and $h$ is its $i^*$-th coordinate
point. Hence $d(x_T,h)^2=|c+Tu|^2+1-2(c_{i^*}+Tu_{i^*})$, which is the claim. For $y_T$ replace $u$ by
$-u$.
\end{proof}

\subsection{The placement}\label{sec:placement}
Put $V=\F_q^2\oplus\F_q^m$, $N=m+2$, $M=0\oplus\F_q^m$. The node $a\in\Mat_{m\times2}(\F_q)$ is the plane
$W_a=\{(z,az):z\in\F_q^2\}$, complementary to $M$. The hub of class $x$ and value $y$ is the point
$h^x_y=\langle(x,y)\rangle$, so the hub of node $a$ in class $x$ is $h^x_{ax}\subset W_a$, and $a,a'$ lie
on a common line of class $x$ iff $h^x_{ax}=h^x_{a'x}$. Since $W_a\cap W_{a'}\cong\ker(a-a')$, the planes
of a generic edge are skew, and those of a rank-one edge meet in one point. All hubs are distinct
points, so all hub distances are $\sqrt2$: $\delta_k=\sqrt2$ for every merge.

\emph{Generic edges.} For $e=\{a,a'\}$ with $\rk(a-a')=2$ write $W=W_a$, $W'=W_{a'}$ and choose $R$ with
$V=W\oplus R\oplus W'$. For $T>0$ let $x_e$ be the filtration with $W$ at level $(1+T)/2$, $W\oplus R$ at
level $0$, $V$ at level $-T/2$, and $y_e$ the same with $W'$ in place of $W$.

\begin{proposition}[\stP]\label{prop:generic}
$d(x_e,y_e)=2T+1$, and $d(x_e,h)=d(y_e,h')=\sqrt{T^2+\tfrac12}$ for every point $h\subset W$ and every
point $h'\subset W'$. Hence for every class, $s_e=2T+1-2\sqrt{T^2+1/2}\to1=\delta/\sqrt2$ as $T\to\infty$.
\end{proposition}

\begin{proof}
For a point $h\subset W$ choose a basis $(h,u_1)$ of $W$, a basis $r$ of $R$ and a basis $(u',v')$ of
$W'$. The basis $(h,u_1,r,u',v')$ is adapted to $x_e$ and $y_e$ with level vectors
$(\alpha,\alpha,0,\dots,0,-\beta,-\beta)$ and $(-\beta,-\beta,0,\dots,0,\alpha,\alpha)$,
$\alpha=(1+T)/2$, $\beta=T/2$. Their difference has four entries $\pm(\alpha+\beta)$, so
$d(x_e,y_e)=2(\alpha+\beta)=2T+1$. The point $h$ is the first coordinate point, so
$d(x_e,h)^2=(\alpha-1)^2+\alpha^2+2\beta^2=T^2+\tfrac12$. The same for $y_e$ and $W'$.
\end{proof}

\emph{Rank-one edges.} Let $W\cap W'=\langle i\rangle$ and let $k_1\in P$ be a class with
$(a-a')x_{k_1}\ne0$; for an edge on a common line of class $k_0$ take for $k_1$ the successor of $k_0$
in a fixed cyclic order of $P$. The hubs $h\subset W$ and $h'\subset W'$ of class $k_1$ are distinct from
$i$, so $h,i,h'$ are independent. Complete them to a basis $f=(h,i,r,h')$ and take
$u=(1,0,\dots,0,-1)/\sqrt2$, $x=F^{f,Tu}$, $y=F^{f,-Tu}$. By \cref{lem:farleg} the saving of class
$k_1$ tends to $u_1-u_N=\sqrt2$. Every other point of $W$ has support $\{i\}$ or $\{h,i\}$ and every
other point of $W'$ has support $\{i\}$ or $\{i,h'\}$, so the savings of the other classes tend to
$0$. Each merge receives the full saving $\sqrt2$ from exactly $E_{\mathrm{coll}}/p=E/(g+1)$ edges. On the
remaining rank-one edges we use the same recipe or legs orthogonal to all hubs; their savings are not
counted.

\begin{theorem}[Theorem A; \stP; refereed, \src{review-q1.md}~\S2]\label{thm:A}
For the placement above, as $T\to\infty$,
\[
\min_k\frac{b_k}{E}\ \ge\ \frac{g-q}{\sqrt2(g+1)}+\frac1{g+1}=\frac1{\sqrt2}-\frac{(q+1)/\sqrt2-1}{g+1}.
\]
Together with \cref{cor:cap},
\[
\frac1{\sqrt2}-\frac{(q+1)/\sqrt2-1}{g+1}\ \le\ \Fcat\bigl(\cN(q,m,P)\bigr)\ \le\ \sqrt{\frac g{2(g+1)}},
\]
for every prime power $q$, every $m\ge2$ and every $P$ with $p\ge2$. On every balanced set $C$ of
generic edges with the induced merges, $\Fcat(C)=1/\sqrt2$.
\end{theorem}

\begin{proof}
$b_k=\sum_es^k_e/\sqrt2$. Generic edges contribute $1/\sqrt2$ each in the limit
(\cref{prop:generic}), the $E/(g+1)$ rank-one edges assigned to $k$ contribute $1$ each, and all
other savings are $\ge0$ up to $O(1/T)$. Insert~\eqref{eq:counts}. For a balanced set $C$ of generic
edges the placement restricted to $C$ gives $1/\sqrt2$, and \cref{thm:cap2} applied to every pair of
merges gives the upper bound $1/\sqrt2$.
\end{proof}

The configurations $X_q(P)$ and $E_D$ are balanced (\cref{lem:XD}), so $\Fcat(X_4)=\Fcat(E_D)=1/\sqrt2$
exactly. A single wire is fully crossing but not balanced; its CAT(0) cap is $1$, so the equality
``value $=$ cap'' needs balance.

\begin{remark}[mechanism]
In the apartment of the basis of \cref{prop:generic} the legs lie on a line through the points
$z(a)$ and $z(a')$, ``$W$ at level $\tfrac12$'' and ``$W'$ at level $\tfrac12$'', which are at distance
$1=\delta/\sqrt2$. The $q+1$ hubs of $a$ are pairwise at distance $\sqrt2$ with common midpoint $z(a)$,
and $x_e$ is equidistant from all of them. A tree cannot do this: a far point is equidistant from
the tips of a tripod only above its centre. The flats of $\cC_{m+2}(q)$ have dimension $m+2$; for
$m=1$ any two node planes meet, and then $\sum_ks^k_e\le\sqrt2$ for every edge
(\src{q1-notes.md} Lemma~6.1, \stR).
\end{remark}

\begin{corollary}[quotients; \stR, \src{q1-notes.md} Cor.~2.2]\label{cor:quot}
If a group $\Gamma$ maps onto $\F_q^m$, the 3-net of $\Gamma$ (nodes $(a,b)$, classes $a$, $b$, $ab$) has
$\Fcat\ge2^{-1/2}(1-q^{-m})(1-q^{1-m})$. This tends to $1/\sqrt2$ for $\Z_n^m$ ($m\to\infty$), for
$\Z_p\times\Z_p$ ($p\to\infty$) and for groups with elementary abelian quotients of rank $\ge2$ and
unbounded order. It gives nothing for cyclic groups.
\end{corollary}

\begin{proof}[Sketch]
Pull the hubs back along $\pi$ (hubs with equal images coincide, so $\delta_k=\sqrt2$); use the legs of
\cref{prop:generic} for edges whose images are generic and legs orthogonal to all hubs for the rest.
Count the generic fraction.
\end{proof}

\emph{Numerical checks} (\stN). The referee's own filtration arithmetic over $\F_2,\F_3,\F_4$
(\src{review-q1.md}~\S2, scratch scripts listed there) gives, at $T=10^4$, saving $0.999950$ on every
generic edge and $\min_kb_k/E=0.48279,0.58247,0.64111$ for $(q,m)=(2,2),(2,3),(2,4)$ against the formula
$0.48284,0.58252,0.64115$ and the caps $0.63246,0.66667,0.68599$; the edge counts
match~\eqref{eq:counts}. An independent Rust implementation reproduces $0.482843$ and $0.582516$
(\src{s3-experiments.md}~\S5c).

% =====================================================================
\section{The price barrier}\label{sec:price}
% =====================================================================

\subsection{Theorem B}

\begin{theorem}[Theorem B; \stR, \src{q1-notes.md}~\S3, \src{review-q1.md}~\S3; \cond{the failure of contracting over $k$ and on \cref{imp:AN}(a),(d)}]\label{thm:B}
Suppose the contracting hypothesis fails over $k$. Let $q=2^s$ and let $\mathcal P$ be a finite subset of
$\cC_N(q)$. Then the restriction of $d$ to $\mathcal P$ is a price over $k$.
\end{theorem}

\begin{proof}[Sketch]
\emph{Step 0.} Replacing every level $s$ by $\tfrac12+\eps s$ multiplies $d$ by $\eps$, so we may assume
all levels lie in $[0,1]$. \emph{Step 1.} For each pair of points of $\mathcal P$ fix a common adapted
basis (\cref{lem:filt}(a)); let $\mathcal B$ be these bases and $\cF$ the finite set of all their coordinate
subspaces. \emph{Step 2: $(\cF,\tau)$ is a price, $\tau(U,U')^2=\dim(U+U')-\dim(U\cap U')$.} Restriction
of scalars to $\F_2$ multiplies $\tau^2$ by $s$. The map $U\mapsto U\oplus U^0\subset\F_2^n\oplus(\F_2^n)^*$
($n=sN$, $U^0$ the annihilator) sends subspaces to $n$-dimensional subspaces with
$\dim((U\oplus U^0)+(U'\oplus U'^0))-n=\tau(U,U')^2$. Over $\F_Q$, $Q=2^t>|\cF|$, these spaces have a
common complement, so they are graphs of maps $a_U$ with $\dim(\cdot+\cdot)-n=\rk(a_U-a_{U'})$, and
the regular representation of $\F_Q$ multiplies ranks by $t$. This gives $\iota:\cF\to\Mat_\mu(\F_2)$,
$\mu\ge3$, with $\tau^2=\rk(\iota U-\iota U')/(st)$, and \cref{imp:AN}(d) makes $\tau$ a price.
\emph{Step 3.} By \cref{imp:AN}(a) there are a price space $Y$ and an isometric $j:\cF\to Y$.
\emph{Step 4.} For $f\in\mathcal B$ the $2^N$ coordinate subspaces carry the metric of the vertices of
$[0,1]^N$, and filling one coordinate at a time (segment filling from the flat-rectangle property,
\src{notes.md} Lemma~4.1, \stR) gives an isometric embedding $\Phi_f:[0,1]^N\to Y$ with
$\Phi_f(1_S)=j(\langle f_i:i\in S\rangle)$. \emph{Step 5.} A point of $\mathcal P$ with chain
$0\subsetneq V_1\subsetneq\dots\subsetneq V_k=V$ and levels $1\ge s_1>\dots>s_k\ge0$ has, in every adapted
$f$, the level vector $(1-s_1)1_0+\sum_j(s_j-s_{j+1})1_{V_j}$ ($s_{k+1}=0$), a convex combination of the
chain vertices with coefficients independent of $f$. Two isometric embeddings of a Euclidean simplex
into a CAT(0) space that agree on the vertices agree everywhere (induction on the dimension, by
uniqueness of geodesics), so $\Phi(p):=\Phi_f(x)$ does not depend on $f$. \emph{Step 6.} For $p,p'\in\mathcal
P$ the basis chosen for the pair is adapted to both, so $\dd{\Phi(p)}{\Phi(p')}=d(p,p')$, and the pull-back
of a price is a price.
\end{proof}

The referee checked Step~2 by building $\iota$ explicitly in two cases (all pairs, no failure), and
noted that restriction of scalars $\cC_N(2^s)\to\cC_{sN}(2)$ is $\sqrt s$ times an isometry on all
filtrations, so the case $q=2$ already implies the others. For odd $q$ no route from $\F_q$-subspace
lattices to $\F_2$-rank metrics is known, and the price statement is open. With all bases in place of
$\mathcal B$ the same argument shows that the closed convex hull of $(\mathrm{Sub}(\F_q^N),\tau)$ in any
CAT(0) space containing an isometric copy of it is the union of the cubes of all bases
(\src{q1-notes.md} Cor.~3.2, \stR).

\subsection{Consequences for one level}

\begin{theorem}[\stR, \src{q1-notes.md}~\S4.1; \cond{the hypotheses of \cref{thm:B}}]\label{thm:C}
Let $q=2^s$, let $C$ be a finite set of generic edges of some $\cN(q,m,P)$ with the induced merges, and
suppose contracting fails over $k$. For every $\eps>0$ and $\tau(\lambda)=o(\lambda)$ there is $\lambda_0$
such that for $\lambda\ge\lambda_0$ every identity $\sum_{i\in I}A_i=\zeta f_C$ over $k$ with
$|I|\le2^{\tau(\lambda)}$ has a piece with $c_k(A_i)\ge(1/\sqrt2-\eps)|C|\lambda$ for all merges $k\in P$ at
once. If $C$ is balanced, no price gives more than $|C|\lambda/\sqrt2$.
\end{theorem}

\begin{proof}
The placement of \cref{sec:placement} restricted to $C$, at a fixed $T=T(\eps)$, is a finite subset of
$\cC_{m+2}(q)$, hence a price by \cref{thm:B}. Apply \cref{prop:price}. The upper bound is
\cref{thm:cap2} for each pair of merges.
\end{proof}

\begin{corollary}[per-level form; \stR; \cond{the hypotheses of \cref{thm:B}}]\label{cor:perlevel}
If contracting fails over $k$, then for every $g=2^m\ge4$, $\eps>0$ and $\tau=o(\lambda)$, every identity over
$k$ for the level tensor of the 3-net of $\F_2^m$ with at most $2^{\tau(\lambda)}$ terms has a piece with
$c_k\ge(1/\sqrt2-1.13/(g+1)-\eps)E\lambda$ for $k=1,2,3$ at once, for $\lambda\ge\lambda_0(g,\eps,\tau)$. For
nets over $\F_{2^s}$ the same holds with $(q+1)/\sqrt2-1$ in place of $1.13$.
\end{corollary}

The constant is $3/\sqrt2-1=1.1213$. The dependence of $\lambda_0$ on $g$ is not controlled: the proof
of \cref{imp:AN}(d) takes a limit of prices, so no rate in $\lambda$ comes out. This is why the next
subsection works with a fixed configuration.

\subsection{A fixed configuration and the positive target}

\begin{lemma}[\stP]\label{lem:XD}
(a) Let $\{a,a'\}$ be a generic edge of $\cN(q,m,P)$ and $S=\operatorname{im}(a-a')$. The coset
$a+\mathrm{Hom}(\F_q^2,S)$ is a sub-net isomorphic to $\cN(q,2,P)$, with lines restricting to lines of the
same class, and it is the only such coset containing the edge. Hence the generic edges of
$\cN(q,m,P)$ are partitioned into copies of $X_q(P)$, with remainder $(q+1)/(g+1)$.
(b) $X_q(P)$ is balanced and fully crossing. (c) $E_D$ is balanced and fully crossing for every
$D\in\GL_2(\F_2)$.
\end{lemma}

\begin{proof}
(a) If a coset $a''+\mathrm{Hom}(\F_q^2,S')$ contains $a$ and $a'$, then $\operatorname{im}(a-a')\subseteq S'$, so
$S'=S$. Inside the coset, line membership and rank are computed as in $\cN(q,2,P)$ after choosing
coordinates on $S$. (b) Let $G,G'$ be lines of classes $x\ne x'$ meeting in $n_0$. Write
$y=n_0+u\varphi_x^{\top}$ and $z=n_0+u'\varphi_{x'}^{\top}$ with $\varphi_x$ vanishing on $x$; then $y-z$ is
invertible iff $u,u'$ are independent, which gives $|\GL_2(\F_q)|=\deg n_0$ ordered pairs.
(c) For $y=n_0+u\varphi_x^\top\in G$ the condition $y+D\in G'$ reads $u\,\varphi_x(x')=Dx'$, which has
exactly one solution; $\deg n_0=1$. Full crossing holds because $Dx\ne0$ for $x\ne0$.
\end{proof}

The referee confirmed the partition numerically for $g=4,8,16$ ($1$, $28$, $560$ copies of $48$ edges,
remainders $0.6,0.3333,0.1765$), the balance of $X_3(P)$ for all $11$ slope sets with $p\ge2$, and the
eight-wire values $b_k=5.656851$ for all six $D$ at $T=10^6$ against $8/\sqrt2=5.656854$
(\src{review-q1.md}~\S\S4--5, \stN).

\begin{corollary}[the positive target; \stR; \cond{\cref{imp:AN}(a),(d)}]\label{cor:target}
If contracting fails over $k$, the price floor of $E_D$ is exactly $8\lambda/\sqrt2$ per merge and that of
$X_4$ is $48\lambda/\sqrt2$. Contrapositively: a sequence $\lambda\to\infty$ of identities
$\sum_iA_i=\zeta f_{E_D}$ over $k$ with $2^{o(\lambda)}$ pieces, each with $\min_kc_k(A_i)\le
(1/\sqrt2-\eta)8\lambda$ for a fixed $\eta>0$, proves the contracting hypothesis over $k$, hence a
power-saving algorithm for 3-SAT (Theorem~R in characteristic $0$, Theorem~C1-AN in characteristic
$p$).
\end{corollary}

This is one fixed finite object with three merges into four groups. For each pair of merges $E_D$ is a
union of four of Anthropic's crossings, so $1/\sqrt2$ is attained for two merges; the content is that
the third merge cannot be used.

\subsection{Engines with per-copy pieces}

\begin{theorem}[Theorem $\mathrm D'$; \stR, \src{q1-notes.md}~\S4.3, \src{review-q1.md}~\S4; \cond{\cref{imp:AN}(a),(d) and Theorem~R / C1-AN}]\label{thm:Dprime}
Let $k$ be a field, $q=2^s$, $P\subseteq\mathbb P^1(\F_q)$ with $p\ge2$ slopes. Consider engines whose
levels are the levels of nets $\cN(q,m_\ell,P)$, with the generic edges partitioned into copies of
$X_q(P)$ as in \cref{lem:XD}(a) (remainder $(q+1)/(q^{m_\ell}+1)$), and whose terms are products of
pieces, one per copy, taken from one identity $\sum_{i\in I}A_i=\zeta f_{X_q(P)}$ over $k$ with
$|I|\le2^{\tau(\lambda)}$, $\tau=o(\lambda)$. If such an engine has log-exponent $\alpha'>1/2$ along
infinitely many $n$, then contracting holds over the algebraic closure $\bar k$, and 3-SAT has a
power-saving algorithm.
\end{theorem}

\begin{proof}[Sketch]
This is the gap theorem of \src{research/ports/c1d1/notes.md}~D1.3 with \cref{thm:C} in place of the flat floor.
An identity over $k$ is one over $\bar k$, and costs over $\bar k$ are at most those over $k$. Put
$\theta=2^{-\alpha'}<1/\sqrt2$. Some level has factor at most $(1-\xi)(1/\sqrt2-\eta)$; then every piece
for $X_q(P)$ at that $\lambda$ is below $(1/\sqrt2-\eta)|X_q(P)|\lambda$ under some merge, since otherwise the
term using an expensive piece on every copy would cost more. The list of configurations has one
member, so this happens for $\lambda\to\infty$, contradicting \cref{thm:C} over $\bar k$ unless
contracting holds there.
\end{proof}

Anthropic's engine is of this per-copy kind (App.~B pieces on each crossing, product terms). The class
is narrower than one would like. It does not contain engines with one identity for the whole net level
and $2^{o(\lambda E)}$ arbitrary terms, nor term sets that are not products over the copies of $X_q$,
such as pieces on whole translation classes of the order-$g$ net. For whole-level identities only
\cref{cor:perlevel} is proved, with $\lambda_0$ depending on $g$, so it gives no statement about towers.
Earlier thresholds $0.7075$ and $0.79$ for whole-level classes remain conjectural for the same reason
(\src{review-sat3.md} item~5).

\subsection{Tree products}\label{sec:trees}
Before \cref{thm:A}, the best price floors came from products of trees, which are prices when
contracting fails (\src{notes.md} Thm.~4.3, \stR, \cond{\cref{imp:AN}(a),(d)}). In a product of nine
tripods, one per node of $\Delta_3$ with the three hubs of that node on the three arms and all other hubs
at the centre, $b=(6,6,6)$, so $\rho_{\Delta_3}\ge2/3$ when contracting fails (\src{notes.md}
Thm.~5.2, \stR). For 3-nets of $\F_2^M$ and $\F_4^M$, a subgroup tripod of index $4$ pulled back from a
quotient gives $\Phi=27/8$ and the limit value $\tfrac12\sqrt{3/2}=0.6124$ (\src{notes.md} Thm.~5.4,
\stR; the transfer holds for $\F_2^M,\F_3^M,\F_4^M$). The exact finite-$g$ values $0.5410,0.5729,0.5915$
at $g=4,8,16$ lie below the flat value for $g\le8$ (\src{review-sat3.md}~4(iii), \stN). For nets of $\F_q^m$
with $m\ge2$ \cref{thm:A} supersedes these values. For cyclic groups and $\Delta_3$ they remain the best known.

% =====================================================================
\section{Theorem W: certified costs of nested crossing covers}\label{sec:thmW}
% =====================================================================

\subsection{The class of covers}\label{sec:class}
The parties are the legs and auxiliary parties created at will. A \emph{wire system} $W$ is a finite list
of wires $(u,v;\beta)$, $\beta$ bits, between parties, with the \emph{leg invariant}: the wires at each
leg carry exactly its $\lambda$ bits, no leg is used as a hub, and no wire at a leg is enlarged (both can be
replaced, at no cost, by routing through an auxiliary party at the leg's group). The start is $W_0=E$.
A move replaces a wire system by one or several, and the tensor it stands for by a cover whose members
are restrictions of the new graph tensors:
\begin{itemize}
\item (M0) split a wire; route a wire through an auxiliary party (Lemma~A.2 of~\cite{AnthropicVP});
increase the bits of a wire between auxiliary parties; identify two auxiliary parties.
\item (M1) \emph{crossing.} Choose two wires $(a,c;\beta_1)$, $(b,d;\beta_2)$ between any parties
(a bond between two hubs is allowed) and a \emph{circle-covering crossing identity} for them: pieces
with certified log-ranks $(\rho_i,\rho_i')$ across $ab|cd$ and $ad|bc$ whose boxes $[0,\rho_i]\times[0,\rho_i']$
cover the arc $\{r^2+r'^2=2(\beta_1^2+\beta_2^2)\}\cap[0,\beta_1+\beta_2]^2$ up to a stated $\eps\ge0$. For
each piece the cover chooses a cut and replaces the two wires by
$(a,o;\beta_1),(b,o;\beta_2),(c,o';\beta_1),(d,o';\beta_2),(o,o';\rho_i)$ (cut $ab|cd$) or by
$(a,h;\beta_1),(d,h;\beta_2),(b,h';\beta_2),(c,h';\beta_1),(h,h';\rho'_i)$ (cut $ad|bc$), with new auxiliary
parties.
\item Sums and regroupings of leaves, certified by direct-sum networks.
\end{itemize}
Instances of (M1) are Anthropic's crossing identity (its Proposition~B.2) with the certificates
$b_c(N)=\log_2((\Theta-N+1)P(\Theta-N))$, $b_b(N)=\log_2((N+1)P(N))$, and Khoruzhii's two-piece integral
identity~\cite{Kho26}, which is a different identity and is inside by circle covering; Sra's nested
identity~\cite[Lemma~4.2]{Sra26} crosses two hub--hub bonds with it. A tree of moves writes
$f=\sum_{\text{leaves}}t_{\text{leaf}}$ with $t_{\text{leaf}}$ a restriction of the graph tensor of
$W_{\text{leaf}}$. A leaf is evaluated under a merge $k$ at the \emph{certified cost}
\begin{equation}\label{eq:certcost}
\cost_k(W)=\min_{\pi:\mathrm{Aux}\to\text{groups of }k}\ \sum_{(u,v;\beta)\in W}\beta\,[\pi(u)\ne\pi(v)],
\end{equation}
where legs are mapped to the groups of their nodes. Pushing forward along $\pi$ gives
$c_k(t_{\text{leaf}})\le\cost_k(W_{\text{leaf}})$: the certified cost is an upper bound on the true cost. The
\emph{certified level value} of a cover is $\max_{\text{leaf}}\min_k\cost_k(W_{\text{leaf}})/(\lambda|E|)$.

For a placement in a metric space $Y$ the \emph{Steiner value} of $W$ is
$v(W)=\inf\sum_{(u,v;\beta)}\beta\,\dd{\mathrm{pos}(u)}{\mathrm{pos}(v)}$ over the positions of the auxiliary
parties. Two facts hold in every metric space (\stP): (S1) $v(W_0)=\lambda L$ and (M0) moves do not
decrease $v$ (triangle inequality); (S2) for every leaf and merge, $v(W)\le\lambda H_k+\delta_k\cost_k(W)$:
put each auxiliary party at the hub $o^k_{\pi(y)}$ of an optimal $\pi$ in~\eqref{eq:certcost}; a wire then
has length at most the leg-to-hub distances of its leg ends plus $\delta_k[\pi(u)\ne\pi(v)]$, and the wires
at a leg carry its $\lambda$ bits.

\subsection{Lemma S}

\begin{lemma}[Steiner crossing inequality, Lemma S; \stPL{} (core), refereed, \src{review-thmW.md}~\S2]\label{lem:S}
Let $Y$ be a CAT(0) space, $W$ a wire system with legs at fixed points of $Y$, and $(a,c;\beta_1)$,
$(b,d;\beta_2)$ two of its wires. Let $r_c$ be the least bond $r\ge0$ for which the child with cut $ab|cd$
and bond $r$ has Steiner value at least $v(W)$, and $r_b$ the same for the cut $ad|bc$. Then
\[
r_c^2+r_b^2\le2(\beta_1^2+\beta_2^2),\qquad0\le r_c,r_b\le\beta_1+\beta_2 .
\]
\end{lemma}

\begin{proof}
Write $v=v(W)$ and $\mathrm{rest}(\mathrm{pos})$ for the total length of the wires of $W$ other than the
two. For a configuration $Z^1=(\mathrm{pos}^1,o,o')$ of the first child put
$N_1=v-\mathrm{rest}(\mathrm{pos}^1)-(\beta_1\dd{a^1}o+\beta_2\dd{b^1}o+\beta_1\dd{c^1}{o'}+\beta_2\dd{d^1}{o'})$,
where $a^1=\mathrm{pos}^1(a)$ (the fixed point if $a$ is a leg). The child with bond $r$ has value
$\ge v$ iff $N_1\le r\dd o{o'}$ for all $Z^1$. For $o=o'$ this holds by the triangle inequality, since
$v\le\mathrm{rest}+\beta_1\dd{a^1}{c^1}+\beta_2\dd{b^1}{d^1}$. Hence $r_c=\max(0,\sup N_1/\dd o{o'})$ over
configurations with $o\ne o'$, and the same inequality with $\dd{a^1}{c^1}\le\dd{a^1}o+\dd o{o'}+\dd{o'}{c^1}$
gives $N_1\le(\beta_1+\beta_2)\dd o{o'}$, so $r_c\le\beta_1+\beta_2$. Define $N_2$ and $r_b$ in the same way
for configurations $Z^2=(\mathrm{pos}^2,h,h')$ of the second child, with
$N_2=v-\mathrm{rest}(\mathrm{pos}^2)-(\beta_1\dd{a^2}h+\beta_2\dd{d^2}h+\beta_2\dd{b^2}{h'}+\beta_1\dd{c^2}{h'})$.

Fix $Z^1$, $Z^2$ with $o\ne o'$, $h\ne h'$, and $a_1,a_2\ge0$ not both $0$. Put $\lambda_1=a_1/\dd o{o'}$,
$\lambda_2=a_2/\dd h{h'}$, $\Lambda=\lambda_1+\lambda_2$, $t=\lambda_2/\Lambda$. For every auxiliary party $y$ of
$W$ let $\mathrm{pos}^t(y)=\gamma_t(\mathrm{pos}^1(y),\mathrm{pos}^2(y))$, and write $a^t,b^t,c^t,d^t$ for the
resulting positions of the four ends. Put $z_a=\gamma_t(o,h)$, $z_c=\gamma_t(o',h')$, $z_b=\gamma_t(o,h')$,
$z_d=\gamma_t(o',h)$; these pair each end with its hub in the two children.
\begin{enumerate}[label=(\roman*)]
\item By convexity of the metric along pairs of geodesics, $\mathrm{rest}(\mathrm{pos}^t)\le(1-t)\mathrm{rest}(\mathrm{pos}^1)
+t\,\mathrm{rest}(\mathrm{pos}^2)$.
\item $\mathrm{pos}^t$ is a configuration of $W$, so $v\le\mathrm{rest}(\mathrm{pos}^t)+\beta_1\dd{a^t}{c^t}+\beta_2\dd{b^t}{d^t}$.
\item By the same convexity, $(1-t)\dd{a^1}o+t\dd{a^2}h\ge\dd{a^t}{z_a}$, and likewise for $b$ (with $o,h'$),
$c$ (with $o',h'$) and $d$ (with $o',h$).
\item $\dd{a^t}{c^t}\le\dd{a^t}{z_a}+\dd{z_a}{z_c}+\dd{z_c}{c^t}$ and
$\dd{b^t}{d^t}\le\dd{b^t}{z_b}+\dd{z_b}{z_d}+\dd{z_d}{d^t}$.
\end{enumerate}
Combining (i)--(iv),
\[
(1-t)N_1+tN_2\le\beta_1\dd{z_a}{z_c}+\beta_2\dd{z_b}{z_d}\le\sqrt{\beta_1^2+\beta_2^2}\,
\bigl(\dd{z_a}{z_c}^2+\dd{z_b}{z_d}^2\bigr)^{1/2}.
\]
By \cref{prop:bimedian} for $o,h,o',h'$ the last bracket is at most $2(1-t)^2\dd o{o'}^2+2t^2\dd h{h'}^2$.
Multiply by $\Lambda$; since $\Lambda(1-t)\dd o{o'}=a_1$ and $\Lambda t\dd h{h'}=a_2$,
\[
a_1\frac{N_1}{\dd o{o'}}+a_2\frac{N_2}{\dd h{h'}}=\Lambda\bigl((1-t)N_1+tN_2\bigr)\le\sqrt{2(\beta_1^2+\beta_2^2)}\,
\sqrt{a_1^2+a_2^2}.
\]
Taking suprema, $a_1r_c+a_2r_b\le\sqrt{2(\beta_1^2+\beta_2^2)}\,|a|$ for all $a\ge0$ (a coordinate that is
$0$ by the $\max$ is covered by the case $a_1=0$ or $a_2=0$). With $a=(r_c,r_b)$ this is the claim.
\end{proof}

The step without an analogue for fixed hub positions is (ii): instead of looking for one bond that is
good at the parent's optimum, the proof interpolates two bad child configurations, and the interpolation
is again a configuration of the parent. The referee's random tests (\src{review-thmW.md}~\S2, \stN) found
no violation in more than $650$ instances in $\R^2$, $\R^3$ and metric trees, including the degenerate
shapes, and found violations of size $0.59$ in the $\ell^1$ and $\ell^\infty$ planes, which are not CAT(0).

\subsection{The zero-loss lemma}
Let $P(E)=\sum_{m\le E/2}p(m)$, $f(E)=\log_2((E+1)P(E))$, and for classes $\Sigma_1,\Sigma_2$ of symbols of
energies $e_i=2m_i$ put $\beta_i=\log_2|\Sigma_i|$ and $\Theta=2(e_1+e_2)$, so that the certificates of
Anthropic's crossing identity are $b_b(N)=f(N)$, $b_c(N)=f(\Theta-N)$.

\begin{lemma}[zero loss; \stR, \src{attain-notes.md} Lemma~5.3, \src{review-lemma53.md}; uses the quoted bound of~\cite{BO16}]\label{lem:zeroloss}
For every $(r,r')$ with $r^2+r'^2=2(\beta_1^2+\beta_2^2)$ and $0\le r,r'\le\beta_1+\beta_2$ there is
$N\in\{0,\dots,\Theta\}$ with $f(N)\ge r'$ and $f(\Theta-N)\ge r$.
\end{lemma}

The written proof has three parts: an analytic argument for $\max(m_1,m_2)\ge31$ from the two-sided bound
$(\sqrt3/(12n))(1\mp n^{-1/2})e^{\mu(n)}\lessgtr p(n)$, $\mu(n)=(\pi/6)\sqrt{24n-1}$, which goes back to
Lehmer and is quoted in the form given by Bessenrodt and Ono~\cite{BO16} (checked against that source);
an exact computation of $(E+1)P(E)\ge\exp(\kappa\sqrt{E+1/2})$, $\kappa=\pi/\sqrt3$, for $44\le E<800$; and
an exact check of all $7260$ pairs $1\le m_1\le m_2\le120$ (scripts \src{attain/zeroloss.py}). When the other
reports were written this lemma had not been refereed. A separate AI referee has since checked it
(\src{review-lemma53.md}) and found it correct, with presentational repairs: state once that the arc statement
implies the same statement for every point of the disc $r^2+r'^2\le2(\beta_1^2+\beta_2^2)$ in the box
$[0,\beta_1+\beta_2]^2$ (such a point is dominated coordinatewise by a point of the arc), allow energy $0$, and
distinguish the constant $\kappa=\pi/\sqrt3$ from the constant of Anthropic's Lemma~B.1. That referee recomputed the
finite parts in interval arithmetic (all $7380$ pairs $0\le m_1\le m_2\le120$, no failure) and took the quoted
bound as given. Earlier, the referee of Theorem~W had computed the exact discretisation loss with the true
partition counts and the true largest classes: zero for all $11\,175$ pairs of even energies up to $300$ and on a
grid up to $1200$ (\src{review-thmW.md}~\S3, \stN). Without \cref{lem:zeroloss} the refereed
discretisation lemma (\src{attain-notes.md} Lemma~5.2) gives a loss $\eps(\beta_1,\beta_2)=o(\beta_1+\beta_2)$
per crossing, but at least $\kappa/\ln2\approx2.6$ bits.

\subsection{Theorem W}

\begin{theorem}[Theorem W; \stR, \src{attain-notes.md}~\S5.3, \src{review-thmW.md}; exact form uses \cref{lem:zeroloss}]\label{thm:W}
Let a configuration be placed in a CAT(0) space. For every crossing-generated cover there is a leaf with,
for every merge $k$ at once,
\[
\cost_k(W_{\text{leaf}})\ \ge\ \Bigl[\lambda(L-H_k)-D\sum_{s}\eps_s\Bigr]\Big/\delta_k,
\]
the sum over the crossings on the branch of the leaf, $\eps_s$ the covering defect of the identity used
at $s$ and $D$ the diameter of the leg points. For covers whose crossings use Anthropic's
identity, \cref{lem:zeroloss} gives $\eps_s=0$, hence $\cost_k(W_{\text{leaf}})\ge\lambda b_k$ exactly, at every
depth, for any number of crossings and any number of terms. Without \cref{lem:zeroloss} the bound is
non-trivial for branches with $o(\lambda)$ crossings.
\end{theorem}

\begin{proof}[Sketch]
The adversary follows one branch. By (S1) it starts at $v=\lambda L$. At a crossing it answers with a piece
whose box contains $(r_c-\eps,r_b-\eps)$, which exists by \cref{lem:S} and circle covering (the point $(r_c,r_b)$ of the disc is dominated by a point of the arc); with $\eps=0$
neither child is below the parent, whichever cut is chosen. With $\eps>0$, pass to the completion and
project near-optimal positions to the closed convex hull of the legs, which has diameter $\le D$; a bond
smaller by $\eps$ lowers $v$ by at most $\eps D$. The other realisations in the class and the (M0) moves
never lower $v$. At the leaf apply (S2).
\end{proof}

\emph{What Theorem W bounds.} It bounds $\cost_k(W_{\text{leaf}})$, the cost of the realisation that the leaf
inherits from its derivation: rank certificates of the crossing identities, routing networks,
identification and nesting. That is an upper bound on the true cost $c_k(t_{\text{leaf}})$, and Theorem~W says
nothing about $c_k(t_{\text{leaf}})$. It is a barrier for a proof method: an analysis whose cost bounds come only
from these certificates cannot certify a level value below $\Fcat$. The gap is real. For the classes
$(2,12),(2,12)$ ($\Theta=48$) the piece $\Psi_4$ has true rank $7$ across $ad|bc$ against the certificate
$20$, and $\Psi_{24}$ has rank $15$ across the third cut $ac|bd$ against the trivial $25$ (referee's exact
computation, \src{review-thmW.md}~\S1, \stN). The true cost minimises over all networks on the groups, not
only the inherited one. True costs are bounded below only through prices, that is, by
\cref{sec:price}, where Theorem~W is not needed. The class excludes rank-one decompositions, whose boxes
$\{(0,0)\}$ do not cover the arc: circle covering is the operative hypothesis, and the bound is tight on
Anthropic's rectangle level.

\begin{corollary}[\stR; uses \cref{lem:zeroloss}]\label{cor:Wnets}
On the $p$-net of $\F_q^m$, $m\ge2$, every crossing-generated cover has certified level value at least
$1/\sqrt2-((q+1)/\sqrt2-1)/(g+1)$, for every prime power $q$, including odd $q$, with no hypothesis on
$k$. On $\Delta_3$ it is at least $2/3$ (nine tripods). Level-1 covers belong to the class, so the best
certified value is at most the two-merge rectangle value, which tends to $1/\sqrt2$; for the 3-net of $\F_2^2$ the window
is $[0.4828,0.6243]$.
\end{corollary}

\begin{proof}
At each fixed $T$ the placement of \cref{sec:placement} consists of points of the CAT(0) space $\cC_{m+2}(q)$
(\cref{fact:cat}). Apply \cref{thm:W} and let $T\to\infty$ (\cref{thm:A}). The tripod product of
\cref{sec:trees} is CAT(0), and no price hypothesis is used here.
\end{proof}

\subsection{Level one}
At level one a cover crosses pairs of edge wires with arbitrary splits $\beta_1:\beta_2$ and realises
each piece with hubs at groups; a term is evaluated under one merge, and the value is
$\min_wh(w)$ over weight vectors $w$ in the simplex, $h(w)$ the value of a linear programme over designs at weights $w$
(\src{attain-notes.md}~\S2).

\begin{proposition}[\stR, \src{attain-notes.md}~\S2.2 with the referee's correction]\label{prop:rect}
On a 3-net the rectangle partner of a generic edge for two merges has its four nodes in four different
groups of the third merge (three if the partner is of the third class). At the adversary's angle
$\theta=\pi/4$ and equal splits the certified cost vector of a rectangle piece is
$(1/\sqrt2,1/\sqrt2,3/2)$ per bit. Hence $\langle w,c\rangle\ge1/\sqrt2$ for all weights, with equality only
when the weight of the third merge is $0$.
\end{proposition}

\begin{table}[h]
\small
\begin{tabular}{@{}lrcccc@{}}
\toprule
level & $|E|$ & $w=(\tfrac12,\tfrac12,0)$ & two-merge formula & $(0.45,0.45,0.1)$ & uniform $w$\\
\midrule
$K_4$, three matchings & 6 & 0.569036 & 0.569036 & 0.609666 & 0.666667\\
AG$(2,3)$ & 36 & 0.603553 & 0.603553 & 0.643198 & 0.695399\\
net of $\Z_2^2$ & 120 & 0.624264 & 0.624264 & 0.680358 & 0.775805\\
net of $\Z_4$ & 120 & 0.624264 & 0.624264 & 0.674678 & 0.752489\\
net of $\Z_5$ & 300 & 0.638071 & 0.638071 & 0.690931 & 0.773945\\
$\Delta_3$ alone & 9 & 0.9085 & -- & 0.9183 & 0.9395\\
\bottomrule
\end{tabular}
\caption{Level-one certified game values $h(w)/|E|$ (\stN; \texttt{\small attain/game1.py}; the referee
reproduced the $K_4$ and AG$(2,3)$ entries with an independent LP).}\label{tab:level1}
\end{table}

With two merges the game returns exactly the rectangle value for every $g$, and $h$ increases away from
the two-merge corners in every case computed (\cref{tab:level1}).

% =====================================================================
\section{Theorem T: algebraic clique engines}\label{sec:thmT}
% =====================================================================

Let $T_K(\lambda)=\bigotimes_{u<v}\langle2^\lambda\rangle_{uv}$ be the wire tensor of the $K$-clique, viewed on
$K$ rows of side $2^{\lambda(K-1)}$, and $f_{K,\lambda}=\sum_jT_{j_1\dots j_K}X_{1,j_1}\cdots X_{K,j_K}$ its
polynomial, one variable array per row. This is the value of Anthropic's $(K,\lambda)$-clique as a
polynomial in the node tables (Def.~3 of~\cite{AnthropicETH}). A \emph{port} property is defined in
OpenAI's Definition~11.1 of~\cite{OpenAIEG}: a field has ports if bounded-degree regular multigraphs of
unbounded order satisfy a uniform Poincar\'e inequality for squared distances in all spectral spaces over
the field; this depends only on the characteristic.

\begin{imported}[OpenAI~\cite{OpenAIHC}, Thm.~3.1; unrefereed, not formalised]\label{imp:HC}
There are an absolute integer $a$ and constants $c,C>0$ such that for every sufficiently large even $d$,
with $m=\lceil\log_2d\rceil$ and $w=\lceil adm^{3/2}\rceil$, there is a $\{0,1\}$-tensor $T_d$ on $d$ rows of
common side $q_0$, $w\le\log q_0\le Cw$, such that over every field whose characteristic has no ports the
degree-component border size of $f_{T_d^{\otimes h}}$ is at least $\exp(chdw/\sqrt m)$ for all $h\ge1$. Here
$T^{\otimes h}$ is the grouped power, with words of length $h$ as row indices.
\end{imported}

\begin{imported}[OpenAI~\cite{OpenAIIM}, Thm.~7.1; unrefereed]\label{imp:IM}
If characteristic zero admits ports, then for some $\delta>0$ a deterministic algorithm decides 3-SAT on $v$
variables in $2^{O(v^{1-\delta})}$ bit operations times a polynomial in the input length.
\end{imported}

\subsection{Pricing the clique with one site per row fails}

\begin{lemma}[Lemma N; \stPL, \lean{LemmaN.lemmaN}]\label{lem:N}
Let $x_1,\dots,x_K$ be points of a pseudometric space, $K$ even, $K\ge4$. There is a partition
$[K]=A\sqcup B$ with $|A|=|B|=K/2$ such that for all points $y_0,y_1$
\[
(K-1)\Bigl[\sum_{i\in A}d(x_i,y_0)+\sum_{i\in B}d(x_i,y_1)\Bigr]\ \ge\ \sum_{u<v}d(x_u,x_v).
\]
\end{lemma}

\begin{proof}
For every partition into halves and all $y_0,y_1$ the triangle inequality on pairs on the same side gives
$\sum_{\text{same side}}d(x_u,x_v)\le(K/2-1)[\sum_Ad(x_i,y_0)+\sum_Bd(x_i,y_1)]$, since each point lies in
$K/2-1$ such pairs. For a uniformly random partition into halves a pair is on the same side with probability
$(K/2-1)/(K-1)$, so some partition has $\sum_{\text{same side}}d\ge\frac{K/2-1}{K-1}\sum_{u<v}d$. Combine and
divide by $K/2-1>0$.
\end{proof}

The statement fails at $K=2$ (take $y_0=x_1$, $y_1=x_2$; Lean \lean{LemmaN.lemmaN_false_K2}). Equality holds for a $K$-pod with unit legs and
both suppliers at the centre. In OpenAI's comparison a term of rank $R$ on a balanced cut, with the rows
as sites and suppliers $y_0,y_1$, is charged $\lambda(K-1)\sum_id(x_i,y_{b(i)})+d(y_0,y_1)\log R$, and the
target $T_K$ costs $\lambda\sum_{u<v}d(x_u,x_v)$. By \cref{lem:N} the first is at least the second for the
cut of the lemma, for every choice of suppliers, in every spectral space and with or without ports
(\stP). So pricing $T_K$ with one site per row proves nothing, not even $R>1$. In particular the inequality
$\sum_{u<v}(2-d(x_u,x_v))\le(1-\delta)(K-1)\min_\ell\sigma_\ell$ with $\delta>0$, proposed as a missing lemma in
\src{prior-barriers.md}~\S5.3, fails in every metric space unless $\min_\ell\sigma_\ell\le0$ (\stP; the
literal question of that note has a negative answer).

\subsection{Restriction to OpenAI's tensor}

\begin{lemma}[star lemma; \stP]\label{lem:star}
Let $a\ge1$, $K\ge2$, and $\mathrm{Star}_v(a)=\bigotimes_{u\ne v}\langle2^a\rangle_{vu}$.
(i) $T_K(2a)=\bigotimes_v\mathrm{Star}_v(a)$, grouped by rows.
(ii) Every tensor $A$ on $K$ rows whose rows $u\ne v$ have side at most $2^a$ is a row-wise restriction of
$\mathrm{Star}_v(a)$.
(iii) For tensors $A^{(1)},\dots,A^{(K)}$ as in (ii), $A^{(v)}$ with centre $v$, $T_K(2a)$ restricts row-wise to
$A^{(1)}\otimes\cdots\otimes A^{(K)}$.
(iv) For every $K$-row tensor $T$ of common side $q$ and every $h_0$ with $q^{h_0}\le2^a$, $T_K(\lambda)$
restricts row-wise to $T^{\otimes Kh_0}$ for every $\lambda\ge2a$.
\end{lemma}

\begin{proof}
(i) Each edge $\{u,v\}$ has $2a$ bits; give $a$ of them to the star of $u$ and $a$ to the star of $v$.
(ii) At row $v$ apply $M_v:|(j_u)_{u\ne v}\rangle\mapsto\sum_iA_{i,(j_u)}|i\rangle$ and at the other rows the
identity or a coordinate projection; then $(M_v\otimes\mathrm{id})\mathrm{Star}_v(a)=\sum_j(\sum_iA_{i,j}|i\rangle)
\otimes\bigotimes_u|j_u\rangle=A$. (iii) Tensor the maps of (ii) over $v$: the local map at row $u$ is $M_u$ on
its own star factor tensored with coordinate projections on the $K-1$ leaf factors. (iv) Take
$A^{(v)}=T^{\otimes h_0}$ for every $v$; the result is $T^{\otimes Kh_0}$ up to a row-local permutation of
letters. For $\lambda>2a$ first project the extra bits to $0$.
\end{proof}

A row of $T_K(2a)$ has $2a(K-1)$ bits and the hosted power at most $Ka$, an overhead $2(K-1)/K<2$. The
referee checked the lemma by explicit contraction for $K=3,4$ (\src{review-q2.md} item~3, \stN).

\begin{lemma}[Lemma L; \stP]\label{lem:L}
If $U=(\bigotimes_uM_u)T$, then $f_U(X)=f_T(Y)$ with $Y_{u,i}=\sum_j(M_u)_{j,i}X_{u,j}$. The substitution is
homogeneous linear, commutes with taking degree components and is continuous on coefficients, so circuit
size and degree-component border size of $f_U$ are at most those of $f_T$ plus the gates for $Y$.
\end{lemma}

\begin{theorem}[Theorem T; \stR, \src{q2-notes.md}~\S3.3, \src{review-q2.md}; \cond{\cref{imp:HC}}]\label{thm:T}
Let the characteristic of $F$ have no ports. For every sufficiently large even $K$ and every
$\lambda\ge\lambda_0(K)=\Theta(K(\log K)^{3/2})$, every arithmetic circuit over $F$ for $f_{K,\lambda}$, also in
the degree-component border sense, has at least $2^{c_TK^2\lambda/\sqrt{\lceil\log_2K\rceil}}$ gates, with an
absolute $c_T>0$.
\end{theorem}

\begin{proof}[Sketch]
Take $d=K$ and OpenAI's $T_d$, $b_0=\lceil\log_2q_0\rceil$, $a=h_0b_0$ with $h_0=\lfloor\lfloor\lambda/2\rfloor/b_0\rfloor\ge1$.
\Cref{lem:star}(iv) gives $T_K(\lambda)\ge T_d^{\otimes h}$ row-wise with $h=Kh_0$. By \cref{lem:L} a circuit for
$f_{K,\lambda}$ gives one for $f_{T_d^{\otimes h}}$ with at most $K2^{\lambda(K-1/2)}$ extra gates, and
\cref{imp:HC} bounds the latter by $\exp(chKw/\sqrt m)$. The arithmetic $\lfloor\lambda/2\rfloor\ge\lambda/3$,
$\lfloor x/b_0\rfloor\ge x/(2b_0)$ and $w/b_0\ge1/(2C_0)$ gives $hKw/\sqrt m\ge K^2\lambda/(12C_0\sqrt m)$; the
extra gates are negligible once $c_TK/\sqrt m>1$. A second form avoids the extra gates by applying the
row-wise maps inside OpenAI's proof, whose first step uses the target only through a property preserved
by row-wise linear maps.
\end{proof}

An \emph{oblivious algebraic} algorithm is one that, for each $(\lambda,B)$, is a straight-line program in
the table entries (over $\Z$, with constants and divisions by fixed non-zero integers) whose output equals
$f_{K,\lambda}$ as a polynomial identity. Anthropic's engine is of this kind (\src{q2-notes.md}~\S5, re-read by
the referee).

\begin{corollary}[\stR; \cond{\cref{imp:HC,imp:IM}}]\label{cor:T2}
If an oblivious algebraic clique-sum algorithm over $\Q$ has $e(K)\sqrt{\log K}\to0$ along a sequence of even
$K$, for instance $e(K)\le(\log K)^{-\alpha}$ with $\alpha>1/2$ for infinitely many even $K$, then
characteristic $0$ admits ports and 3-SAT on $v$ variables is decidable deterministically in time
$2^{O(v^{1-\delta})}\mathrm{poly}$; so variable-count Weak ETH fails.
\end{corollary}

The referee's sentence on what may be claimed (\src{review-q2.md}, last section) is the following, and we
claim no more. If OpenAI's unrefereed and unformalised HC Thm.~3.1 holds as printed, then over any field
whose characteristic has no ports, every straight-line program that computes Anthropic's clique polynomial
$f_{K,\lambda}$ as a polynomial identity needs $2^{c\cdot K^2\lambda/\sqrt{\log K}}$ gates, for all
sufficiently large even $K$ and all $\lambda\ge C\cdot K(\log K)^{3/2}$, with unspecified absolute constants. If
in addition OpenAI's IM Thm.~7.1 holds (unrefereed; a Lean counterpart \lean{middle_count_bridge} exists in
their release, which we did not build), any such polynomial-identity algorithm over $\Q$ with
$e(K)\sqrt{\log K}\to0$ along even $K$ would give 3-SAT in deterministic time $2^{O(v^{1-\delta})}$ and refute
variable-count Weak ETH. Nothing is claimed in positive characteristic, for algorithms that are correct only
on bounded integer tables, for non-algebraic algorithms, or for wires shorter than $C\cdot K(\log K)^{3/2}$.

In a tower with $\beta=2$ and level factors $\rho_j$, Theorem~T forces the geometric mean of the $\rho_j$ to
be at least $1/\sqrt2$ in the limit. This is the cap of \cref{cor:cap} and the value of the two-merge engine;
Theorem~T makes no per-level price statement and does not answer whether a price reaches $1/\sqrt2$ on
3-nets of cyclic groups. The referee rechecked the arithmetic of the three steps of OpenAI's proof that carry
the bound and found no error, but did not verify its Theorems~5.1 and 6.1, Lemma~7.3 or \S\S9--10
(\src{review-q2.md} item~7).

% =====================================================================
\section{What is open}\label{sec:open}
% =====================================================================

\subsection{Cyclic nets}
\Cref{thm:A} needs an elementary abelian quotient of rank $\ge2$. For cyclic groups, hence for $\Delta_3$
(the generic edges of the $\Z_3$ net) and $K_4$, the realisation by node planes fails.

\begin{proposition}[\stR, corrected after the referee, \src{q1-notes.md} Prop.~6.2, \src{review-q1.md}~6a]\label{prop:Z3}
Let $R_r,C_c,D_d$ ($r,c,d\in\Z_3$) be $s$-dimensional subspaces such that at every node $(r,c)$ the hub
spaces $R_r,C_c,D_{r+c}$ are pairwise transversal and lie in a common $2s$-dimensional space $W_{rc}$. Then all
nine lie in a space of dimension at most $3s$, and any two $W$'s meet in dimension at least $s$.
\end{proposition}

As first stated, without transversality, the proposition was false; the referee's counterexample over $\F_2^4$
has coinciding hubs at some node. Under the hypothesis, cyclic nets have no realisation by subspaces with
skew generic node planes, and point hubs give $\Delta_3$ at most $1/3$. Realisations with coinciding or
non-transversal hubs, hubs at interior points of $\cC_N$, and spaces other than $\cC_N$ are not bounded by
anything here. So $\Fcat(\Delta_3)\in[2/3,1/\sqrt2]$, and for no cyclic net is an upper bound below
$1/\sqrt2$ known.

The evidence on cyclic nets is numerical and was not refereed (\src{q1-cyclic.md},
\src{s3-experiments.md}).
\begin{itemize}
\item A convex relaxation (squared distances on hubs and iterated midpoints, all CN inequalities, convex
1-Lipschitz leg functions, reduced by the $108$ automorphisms of $\Delta_3$) gives $\sum_kb_k\le19.091883
=27/\sqrt2$ at every depth tried, up to $1665$ points: exactly the cap, so it neither proves $6$ nor anything
below $9/\sqrt2$. With every CN inequality forced to equality it still gives the cap unless node barycentres
are included, so the weakness is in the leg functions (\stN). Two lemmas of that note, a flat hub hull
bounds the value by the flat value, and legs may be replaced by convex sets, are proved there and are \stU.
\item Searches over products of trees, cones over $15$ metric graphs including free girth-constrained links,
and filtration spaces $\cC_3(\F_2)$, $\cC_3(\F_3)$ all land on $\Gamma=36$, i.e. $6$ per merge on $\Delta_3$, and
never above (\stN). For full cyclic nets the best tree values are $0.5202$, $0.5519$ at $g=3,5$ against the
flat $0.4979$, $0.5244$, and the subgroup construction of \cref{sec:trees} gives only $0.5857$ for $\Z_g$ in
the limit (\src{review-sat3.md}~4(iii)).
\item The equality analysis of \src{q1-cyclic.md}~\S3 excludes exact equality at both ends of the
symmetric-tripod branch and is undecided in between (\stU).
\end{itemize}

\subsection{Depth-two certified games}
By \cref{thm:W} the certified value of nested covers on a cyclic net can drop at most to $\Fcat$, which is
unknown there. At the two-merge corners no depth helps, because the flat placement already gives the level-one
value. The certified depth-two game was computed on $\Delta_3$ and on the $\Z_5$ net for two strategy types
(\stN; upper bounds for restricted strategy classes, every added strategy verified by an accurate evaluation).
Type~A crosses the bond of a level-one piece with a third wire and may stop or nest adaptively
(\src{attain-notes.md}~\S10, outputs \src{attain/out3/}). Type~B crosses two hub--hub bonds (executor runs,
outputs \src{attain/out4/}; not yet written up in the notes and not refereed).

\begin{table}[h]
\small
\begin{tabular}{@{}llccc@{}}
\toprule
configuration & weights $w$ & level 1 & with type A & with type A and B\\
\midrule
$\Z_5$ net (300 edges) & $(\tfrac12,\tfrac12,0)$ & 0.638071 & 0.638071 & 0.638071\\
& $(0.48,0.48,0.04)$ & 0.659215 & 0.659215 & 0.659215\\
& $(0.45,0.45,0.10)$ & 0.690931 & 0.690931 & 0.690931\\
& $(0.4,0.4,0.2)$ & 0.740858 & 0.740858 & 0.740858\\
& uniform & 0.773799 & 0.771218 & 0.769856\\
$\Delta_3$ (9 edges) & uniform & 0.939002 & 0.929966 & 0.925198\\
& $(\tfrac12,\tfrac12,0)$ & 0.908503 & -- & 0.908503\\
\bottomrule
\end{tabular}
\caption{Certified depth-two game values per edge bit (\stN). Level-one values from the runs of
\texttt{\small attain/out4/} with $11$ split ratios; with $9$ ratios the uniform $\Z_5$ value is $0.773945$.}\label{tab:depth2}
\end{table}

On the $\Z_5$ net the game value stays at the two-merge corner, $0.638071$, for both types; near the corner no
strategy of either type improves on level one, and at uniform weights nesting gains $0.5\%$, far from the drop
below $0.638$ that a gain in the game value would need. On $\Delta_3$ alone, where level one has no rectangles,
nesting is strictly better than level one, and the level-one game value $0.89190$ drops to at most $0.89102$
with type A. The room between the tree record $0.5519$ and $0.6381$ on $\Z_5$ allowed by Theorem~W is not
entered by these strategies. Depth three, third wires depending on the first angle, continuous splits and
$\Z_7$ were not computed.

\subsection{True versus certified costs, and other gaps}
\begin{itemize}
\item \emph{True costs.} The only lower bounds on true costs here are prices (\cref{sec:price}), and they need
the contracting hypothesis to fail. Whether the certificates of Anthropic's crossing identity can be improved
on nested pieces, and whether true ranks below the certificates (\cref{sec:thmW}) can be used, is open.
\item \emph{Odd $q$.} Whether $\sqrt{\rk}$ on $\Mat_\mu(\F_3)$, or $\tau$ on $\mathrm{Sub}(\F_3^N)$, is a price when
contracting fails is open; for nets of $\F_3^m$ only the certified statement holds.
\item \emph{The $O(1/g)$ window.} For $q=2$ the deficit of $\Fcat$ below $1/\sqrt2$ lies between
$0.354/(g+1)$ and $1.121/(g+1)$; within placements with these hubs the rank-one edges cannot do better.
\item \emph{A rate.} A quantitative form of \cref{imp:AN}(d) would make \cref{cor:perlevel} uniform in $g$
and extend \cref{thm:Dprime} to whole-level identities.
\item \emph{Non-algebraic algorithms.} Branching on table values, hashing and instance-adaptive algorithms
are outside every statement of this paper. So are 3-SAT algorithms that do not pass through clique sums, and
short wires $\lambda\ll K(\log K)^{3/2}$ in Theorem~T.
\end{itemize}

% =====================================================================
\section{Limitations}\label{sec:limits}
% =====================================================================

\begin{itemize}
\item No algorithm is given and no unconditional lower bound on true costs is proved. Every barrier concerns a
stated class: per-copy engines (\cref{thm:Dprime}), single net levels (\cref{cor:perlevel}), crossing-generated
covers with certified costs (\cref{thm:W}), or straight-line programs in the long-wire regime (\cref{thm:T}).
\item The price barrier rests on Anthropic's Lean statement \lean{thm_alt_rank_price} and Proposition~C.3, from
a paper that no person has refereed; the consequence for 3-SAT rests on Theorem~R, whose uniform form rests on
an audit by one reader.
\item Theorem T rests on OpenAI's HC Theorem~3.1, which is unrefereed and outside OpenAI's formalised scope;
its constants are unspecified, so no explicit pair $(K,\lambda)$ can be named at which a small circuit would
decide anything.
\item The exactness of Theorem W at all depths rests on \cref{lem:zeroloss}, which was refereed by one AI agent
after the other results and which uses a quoted two-sided bound for the partition function.
\item The hypotheses ``contracting fails'' and ``no ports'' are used in different statements. We do not know
whether either holds; if contracting holds, 3-SAT is power-saving by Theorem~R, and the barriers become moot.
\item All referees were AI agents, the numerical checks ran on one machine, and no person has checked the proofs
or the code.
\end{itemize}

% =====================================================================
\section{Related work}\label{sec:related}
% =====================================================================

Two-merge barriers are now numerous. Majdoddin~\cite{Maj26} defines a rewrite-and-merge class containing
Anthropic's recursion and bounds the rate of two balanced homogeneous merges; his Proposition~4.23 charges a piece
that is local for one merge only under that merge, so it is not a barrier for the game of \cref{sec:thmW}, where
such pieces have certified costs under every merge. De Brito~\cite{deBrito26} shows that the exponent $1/2$ is
optimal for graded pieces of Anthropic's shape, and Fradj~\cite{Fradj26} proves window costs above $1/\sqrt2$ for
code-graded spaces. Sra~\cite{Sra26} proves, conditionally on OpenAI's Euclidean price theorem, the two-merge floor
$1/\sqrt2$ at finite sizes; \cref{thm:cap2} gives the corresponding cap from CAT(0) geometry alone. His nested
identity attains $1/\sqrt k$ on $k$ independent wires with all $2^{k-1}$ transversal cuts; on a net its terms are
certified across two of the cuts only, and it lies in the class of \cref{thm:W}. None of these papers treats three
merges or nets. Our earlier paper~\cite{Tul26barrier} gives barriers for named network models of the counting
engines. OpenAI's \cite{OpenAIEG}, \S11.1, derives ports from a uniform contracting gap on the $\F_2$-Grassmannian;
the filtration space of \cref{sec:thmA} contains the Grassmannian as its $0/1$ points.

% =====================================================================
\section{The Lean statements}\label{sec:lean}
% =====================================================================

The files \src{research/ports/lean/PortsLean/Bimedian.lean} and \src{LemmaS.lean} (Lean v4.36.0 with Mathlib)
work in an arbitrary \lean{PseudoMetricSpace}; Mathlib has no CAT(0) class, so a point $\gamma_t(x,y)$ enters only
through instances of~\eqref{eq:CN}, the predicate \texttt{CN t x y m z}. Formalised:
\lean{PortsLean.Bimedian.bimedian_weighted} (\cref{prop:bimedian} from six CN instances) and
\lean{bimedian_half} ($t=\tfrac12$); \lean{dist_interp_le} (convexity of the metric along two geodesics, from six CN
instances and one auxiliary point, so it is proved and not assumed); \lean{bimedian_of_interp} and
\lean{geoConvex_of_interp} for a global interpolation map \lean{IsCNInterp}; and in inner product spaces
\lean{bimedian_inner_eq} (equality). In \src{LemmaS.lean}: \lean{PortsLean.LemmaS.lemmaS_core} (the inequality
$a_1N_1/\dd o{o'}+a_2N_2/\dd h{h'}\le\sqrt{2(\beta_1^2+\beta_2^2)}\sqrt{a_1^2+a_2^2}$ from hypotheses (i)--(iii) and the
bimedian bound), \lean{lemmaS_interp} (the same with (iii) and the bimedian bound discharged from \lean{IsCNInterp}),
\lean{lemmaS_inner}, \lean{Nc_le} and \lean{Nb_le} ($N\le(\beta_1+\beta_2)\cdot$hub distance),
\lean{circle_of_linear} and \lean{circle_of_sSup} (the circle for the suprema $r_c,r_b$). Not formalised: the
identification of $r_c$ with the least bond, which involves the Steiner value as an infimum over positions;
step~(i) beyond an abstract hypothesis \lean{h_rest}; and step~(ii), the hypothesis \lean{h_parent}. Nothing else
of \cref{sec:thmW} is formalised.

The file \src{research/ports/lean/PortsLean/LemmaN.lean}, added on the same day after the reports above, proves
\cref{lem:N} in the stated form (the cut is chosen before the suppliers) as \lean{LemmaN.lemmaN}, for even $K\ge4$,
by averaging exactly over all balanced cuts (\lean{LemmaN.core}, \lean{LemmaN.sum_cuts_same}), and its failure at $K=2$
as \lean{LemmaN.lemmaN_false_K2}. For the star lemma only the bit counts and the arithmetic of the proof of
\cref{thm:T} are formalised (\lean{LemmaN.star_row_bits}, \lean{LemmaN.star_hosted_bits},
\lean{LemmaN.star_overhead}, \lean{LemmaN.thmT_arith}); the row-wise restrictions are not. Nothing else of this
paper is formalised. We checked on 27 November 2026 that the three files contain no \texttt{sorry}, and
\texttt{\#print axioms} gives \lean{propext}, \lean{Classical.choice} and \lean{Quot.sound} for the fourteen theorems
that we printed: \lean{bimedian_weighted}, \lean{bimedian_half}, \lean{dist_interp_le}, \lean{bimedian_inner_eq},
\lean{lemmaS_core}, \lean{lemmaS_interp}, \lean{lemmaS_inner}, \lean{circle_of_linear}, \lean{circle_of_sSup},
\lean{Nc_le}, \lean{lemmaN}, \lean{lemmaN_false_K2}, \lean{star_overhead} and \lean{thmT_arith}.

% =====================================================================
\section{How this paper was produced}\label{sec:provenance}
% =====================================================================

The mathematics, the proofs, the code, the Lean files and the text of this paper were produced by AI systems
under the author's direction, in one day of work recorded in the repository notes. Each main note was checked by a
separate AI agent acting as referee, with its own scripts; the reports are \src{review-q1.md},
\src{review-sat3.md}, \src{review-thmW.md}, \src{review-q2.md} and \src{review-lemma53.md}, and their corrections are
included above. The
referees were AI agents. No person has checked the proofs, the code or the computations. The notes, reports,
scripts and Lean sources are in the repository~\cite{Repo} at commit \texttt{\commit}:
\begin{center}\footnotesize\expandafter\url\expandafter{\repourl}\end{center}

\bibliographystyle{amsplain}
\bibliography{refs}

\end{document}
